% fntena_logic_copy2.py -- математическая постановка и описание управления.
% Сборка: pdflatex fntena_logic_copy2_math.ru.tex

\documentclass[a4paper,11pt]{article}

\usepackage[T2A]{fontenc}
\usepackage[utf8]{inputenc}
\usepackage[russian]{babel}
\usepackage{amsmath,amssymb}
\usepackage{geometry}
\usepackage{hyperref}
\usepackage{enumitem}
\usepackage{array}
\geometry{left=2.5cm,right=2.5cm,top=2.5cm,bottom=2.5cm}
\setlist{nosep,leftmargin=*}
\hypersetup{colorlinks=true,linkcolor=blue,urlcolor=blue}

\newcommand{\code}[1]{\texttt{\detokenize{#1}}}
\newcommand{\sat}{\operatorname{sat}}
\newcommand{\clip}{\operatorname{clip}}
\newcommand{\sgn}{\operatorname{sgn}}

\title{Математическая постановка сценария\\
\code{fntena_logic_copy2.py}:\\
вертикальная антенна роя БПЛА над якорем}
\author{}
\date{}

\begin{document}
\maketitle

\begin{abstract}
В документе описана математическая модель сценария
\code{Drone\_Swarm\_Simulator\_v2/scenarios/fntena\_logic\_copy2.py}.
Цель сценария --- сформировать вертикальную ``антенну'' из нескольких
мультикоптеров: один аппарат является якорем, остальные должны выстроиться
почти над ним по горизонтали и распределиться по высоте с заданным шагом~$w$.
Движение по $x,y,z$ задаётся локальным законом ближайших соседей, близким к
логике проекта \code{grid\_antenna}. Для $x,y$ сигналы $\sigma_x,\sigma_y$
непрерывно нормируются и переводятся в RC-команды, а для $z$ сигнал
$\sigma_z$ дополнительно проходит через PI-слой, согласующий математический
закон с инерционным поведением \code{POSHOLD}.
\end{abstract}

\tableofcontents
\newpage

%====================================================================
\section{Постановка задачи}

Рассматривается рой из $N$ дронов, где $N \ge 1$. Каждый дрон имеет
целочисленный идентификатор $i \in \{1,\ldots,N\}$. Один дрон с номером
$i=a$ называется \emph{якорем} (в коде параметр \code{--anchor-id}, по
умолчанию $a=1$). Остальные дроны являются агентами вертикальной антенны.

Требуется получить установившееся поведение:
\begin{enumerate}
  \item якорь удерживает заданную высоту $h_a^\star$ и не участвует в
  распределении по высоте;
  \item все агенты сходятся к вертикальной колонне над якорем:
  \[
    x_i \to x_a,\qquad y_i \to y_a,\qquad i \ne a;
  \]
  \item соседние агенты по вертикали расходятся примерно на расстояние
  \[
    w>0,
  \]
  где $w$ --- желаемый шаг антенны в метрах
  (константа \code{W\_SPACING\_M}, аргумент \code{--w});
  \item шаг $w$ должен быть меньше радиуса видимости
  \[
    0 < w < R_{\mathrm{vis}},
  \]
  иначе требуемый сосед не может одновременно быть на расстоянии $w$ и
  оставаться видимым локальному закону.
\end{enumerate}

В текущей версии сценария по умолчанию:
\[
R_{\mathrm{vis}}=2.5\ \text{м},\qquad w=0.8\ \text{м}.
\]

%====================================================================
\section{Система координат и обозначения}

Используется система координат NED (North--East--Down). В коде положение
дрона $i$ записывается как
\[
  p_i = (x_i,y_i,z_i)^\top \in \mathbb{R}^3,
\]
где:
\begin{itemize}
  \item $x_i$ --- координата North, м;
  \item $y_i$ --- координата East, м;
  \item $z_i$ --- координата Down, м; чем меньше $z_i$, тем выше дрон.
\end{itemize}

Физическая высота над точкой старта выражается как
\[
  h_i = -z_i.
\]
Скорость в NED обозначается
\[
  v_i = (\dot x_i,\dot y_i,\dot z_i)^\top
      = (v_{x,i},v_{y,i},v_{z,i})^\top,
\]
где $v_{z,i}>0$ означает движение вниз, а $v_{z,i}<0$ --- движение вверх.

Из-за того, что в SITL каждый экземпляр дрона имеет собственный home-offset
по оси $y$, сценарий переводит локальные позиции в общий кадр:
\[
  y_i^{\mathrm{common}} = y_i^{\mathrm{local}} + (i-1)\Delta_y,
  \qquad \Delta_y = 2.0\ \text{м}.
\]
Далее в формулах под $x_i,y_i,z_i$ понимаются координаты уже в общем кадре.

%====================================================================
\section{Видимые соседи}

Для агента $i$ относительный вектор на дрон $j$ определяется как
\[
  r_{ij}=p_j-p_i=(r_{ij,x},r_{ij,y},r_{ij,z})^\top.
\]

Дрон $j$ считается видимым для $i$, если
\[
  \|r_{ij}\|_2 \le R_{\mathrm{vis}}.
\]
Для неякорных агентов якорный дрон может входить в множество видимых как
обычный сосед, если он находится в радиусе видимости:
\[
  r_{ia}=p_a-p_i,\qquad \|r_{ia}\|_2\le R_{\mathrm{vis}}.
\]

Множество относительных векторов, доступных агенту $i$, обозначим
\[
  \mathcal{N}_i =
  \{r_{ij}\;|\; j\ne i,\ \|r_{ij}\|_2\le R_{\mathrm{vis}}\}
  \cup
  \{r_{ia}\ \text{для неякорного агента, если якорный дрон видим}\}.
\]

%====================================================================
\section{Ближайшие соседи по полуосям}

Для каждой оси $b\in\{x,y,z\}$ агент выделяет ближайший видимый объект в
положительном и отрицательном направлениях. Вводится малый порог
$\varepsilon>0$ (в коде \code{AXIS\_EPS\_M}), чтобы нулевая проекция не
считалась соседом по данной оси.

Для оси $b$:
\[
  d_{b,i}^{+} =
  \min_{r\in\mathcal{N}_i,\ r_b>\varepsilon} r_b,
\]
\[
  d_{b,i}^{-} =
  \min_{r\in\mathcal{N}_i,\ r_b<-\varepsilon} (-r_b).
\]

Здесь:
\begin{itemize}
  \item $d_{b,i}^{+}$ --- расстояние до ближайшего видимого соседа в
  положительном направлении оси $b$;
  \item $d_{b,i}^{-}$ --- расстояние до ближайшего видимого соседа в
  отрицательном направлении оси $b$;
  \item если соседа в соответствующем направлении нет, значение считается
  неопределённым и затем заменяется виртуальным расстоянием.
\end{itemize}

%====================================================================
\section{Заполнение отсутствующих соседей}

Для горизонтальных осей используется анонимное правило соседей. Если по
направлению $x^+$, $x^-$, $y^+$ или $y^-$ ближайшего видимого соседа нет,
эффективное расстояние в этом направлении считается нулевым:
\[
  \tilde d_{x,i}^{+}=0 \text{ при отсутствии соседа в } x^+,\qquad
  \tilde d_{x,i}^{-}=0 \text{ при отсутствии соседа в } x^-,
\]
\[
  \tilde d_{y,i}^{+}=0 \text{ при отсутствии соседа в } y^+,\qquad
  \tilde d_{y,i}^{-}=0 \text{ при отсутствии соседа в } y^-.
\]
Поэтому, если сосед есть только с одной стороны, агент стремится сократить
горизонтальную дистанцию до него до нуля. Если соседей по горизонтальной оси
нет ни с одной стороны, остаётся только скоростной член $\sigma_b=v_b$, то есть
управление гасит собственную скорость.

Для вертикальной оси дополнительно используется направление к якорной точке,
чтобы сохранить растяжение антенны по высоте. Пусть
\[
  q_i = p_a-p_i
\]
--- вектор от агента $i$ к якорю. Для вертикальных направлений строятся знаки:
\[
  s_z^+=\sgn(q_z),\qquad s_z^-=-\sgn(q_z).
\]

Если сосед по вертикальному направлению отсутствует, код подставляет
виртуальное расстояние:
\[
  \tilde d_z =
  \begin{cases}
    R_{\mathrm{vis}}, & \text{если соответствующая компонента } s_z > 0,\\
    w, & \text{если соответствующая компонента } s_z \le 0.
  \end{cases}
\]

Тем самым агент получает шесть эффективных расстояний
\[
  D_i =
  \bigl(
  \tilde d_{x,i}^{+},
  \tilde d_{y,i}^{+},
  \tilde d_{z,i}^{+},
  \tilde d_{x,i}^{-},
  \tilde d_{y,i}^{-},
  \tilde d_{z,i}^{-}
  \bigr).
\]

На практике для управления используются все три пары расстояний:
по ним вычисляются $\sigma_x,\sigma_y,\sigma_z$. Отличие только в
исполнительном слое: $x,y$ переводятся в RC как непрерывное нормированное управление,
а $z$ проходит через PI-регулятор по $\sigma_z$.

%====================================================================
\section{Нелинейный распределённый закон по высоте}

Вводится ограниченная нелинейность
\[
  \xi(g)=\tanh(g).
\]
Она сглаживает влияние расстояний: при больших $g$ значение стремится к~1,
а при малых расстояниях меняется почти линейно.

Для оси $z$ вычисляется скаляр
\[
  \sigma_{z,i}
  =
  v_{z,i}
  -
  \xi(\tilde d_{z,i}^{+})
  +
  \xi(\tilde d_{z,i}^{-}).
\]

Здесь:
\begin{itemize}
  \item $v_{z,i}$ --- текущая вертикальная скорость в NED, положительная вниз;
  \item $\tilde d_{z,i}^{+}$ --- эффективное расстояние до ближайшего объекта
  в направлении $+z$ (ниже агента в NED);
  \item $\tilde d_{z,i}^{-}$ --- эффективное расстояние до ближайшего объекта
  в направлении $-z$ (выше агента в NED);
  \item $\sigma_{z,i}$ --- сигнал рассогласования вертикального распределения.
\end{itemize}

Интуитивно, знак $\sigma_{z,i}$ показывает, в какую сторону необходимо
изменить высоту. В реализации ArduPilot через throttle используется обратная
физическая ось: увеличение throttle заставляет дрон подниматься, то есть
уменьшать координату $z$.

%====================================================================
\section{Горизонтальное управление через соседей}

В соответствии с \code{grid\_antenna} по горизонтальным осям применяется
тот же sigma-закон, что и по вертикали:
\[
  \sigma_{x,i}
  =
  v_{x,i}
  -
  \xi(\tilde d_{x,i}^{+})
  +
  \xi(\tilde d_{x,i}^{-}),
\]
\[
  \sigma_{y,i}
  =
  v_{y,i}
  -
  \xi(\tilde d_{y,i}^{+})
  +
  \xi(\tilde d_{y,i}^{-}).
\]

Здесь $\tilde d_{x,i}^{+},\tilde d_{x,i}^{-},
\tilde d_{y,i}^{+},\tilde d_{y,i}^{-}$ --- эффективные расстояния до
ближайших видимых соседей; если сосед с нужной стороны отсутствует, для
горизонтали подставляется ноль.

Горизонтальное управляющее воздействие строится по той же соседской
ошибке $\sigma$, но вместо дискретного \code{sign}-закона используется
непрерывная нормированная версия, чтобы задействовать промежуточные значения
PWM:
\[
  u_{x,i} =
  \begin{cases}
    0, & |\sigma_{x,i}|<\delta_{xy},\\
    \clip(-\sigma_{x,i},-1,1), & |\sigma_{x,i}|\ge\delta_{xy},
  \end{cases}
\]
\[
  u_{y,i} =
  \begin{cases}
    0, & |\sigma_{y,i}|<\delta_{xy},\\
    \clip(-\sigma_{y,i},-1,1), & |\sigma_{y,i}|\ge\delta_{xy}.
  \end{cases}
\]
Таким образом $u_{x,i},u_{y,i}\in[-1,1]$. В коде $\delta_{xy}$ задаётся
параметром \code{XY\_SIGMA\_DEADBAND}; старые параметры \code{XY\_PWM\_MAX}
и \code{XY\_PWM\_MIN\_STEP} оставлены только для совместимости CLI и не
участвуют в горизонтальном преобразовании.

%====================================================================
\section{PI-регулятор вертикального sigma-сигнала}

Для вертикали используется дискретный PI/PID-регулятор. Для входа $e_k$ на шаге
$k$ с периодом $T_c$:
\[
  I_k=\sat_{[-I_{\max},I_{\max}]}\left(I_{k-1}+e_kT_c\right),
\]
\[
  D_k=\alpha D_{k-1}+(1-\alpha)\frac{e_k-e_{k-1}}{T_c},
\]
\[
  u_k=\sat_{[-U_{\max},U_{\max}]}
  \left(K_P e_k+K_I I_k+K_D D_k\right).
\]

Обозначения:
\begin{itemize}
  \item $K_P,K_I,K_D$ --- пропорциональный, интегральный и дифференциальный
  коэффициенты;
  \item $I_k$ --- интегральная сумма;
  \item $D_k$ --- фильтрованная производная;
  \item $\alpha$ --- коэффициент фильтра производной
  (в коде \code{derivative\_alpha});
  \item $I_{\max}$ --- ограничение интегральной суммы;
  \item $U_{\max}$ --- ограничение выходного PWM-приращения.
\end{itemize}

Для вертикального сигнала используется PI-регулятор по $\sigma_z$:
\[
  u_{z,i}=PI_z\bigl(\operatorname{db}(\sigma_{z,i};\delta_z)\bigr),
\]
где
\[
  K_{P,z}=360,\qquad K_{I,z}=110,\qquad K_{D,z}=0,
\]
\[
  \delta_z=\text{\code{Z\_SIGMA\_DEADBAND}}=0.04,\qquad
  |u_{z,i}|\le \text{\code{Z\_PWM\_MAX}}=220.
\]

%====================================================================
\section{Преобразование в RC PWM}

Нейтральное значение RC-каналов:
\[
  PWM_0 = 1500.
\]

Для нормированного горизонтального управления $u\in[-1,1]$ используется
линейное масштабирование:
\[
  PWM(u)=\clip(1500+400u,1100,1900),
\]
то есть $u=-1$ соответствует $1100$, $u=0$ --- $1500$, $u=1$ --- $1900$.
Тогда горизонтальные команды:
\[
  PWM_{\mathrm{pitch},i}=PWM(-u_{x,i}),
\]
\[
  PWM_{\mathrm{roll},i}=PWM(u_{y,i}).
\]

Вертикальный выход $u_{z,i}$ остаётся PWM-приращением PI-регулятора, поэтому
для throttle используется:
\[
  PWM_{\mathrm{thr},i}
  =
  \sat_{[1100,1900]}(PWM_0+u_{z,i}).
\]

Знак pitch выбран по принятой в сценариях проекта конвенции:
если соседский закон даёт положительное $u_x$ (требуется движение в сторону
увеличения $x$), то в RC override это соответствует отрицательному
нормированному входу канала pitch. Поэтому используется $PWM(-u_{x,i})$.

Для throttle:
\begin{itemize}
  \item $PWM_{\mathrm{thr}}>1500$ --- команда набора высоты
  (уменьшения NED-координаты $z$);
  \item $PWM_{\mathrm{thr}}<1500$ --- команда снижения
  (увеличения NED-координаты $z$).
\end{itemize}

%====================================================================
\section{Управление якорем}

Для якоря $i=a$ горизонтальные команды всегда нейтральны:
\[
  PWM_{\mathrm{roll},a}=PWM_{\mathrm{pitch},a}=PWM_0.
\]

Якорь удерживает высоту
\[
  h_a^\star=\text{\code{TAKEOFF\_ALT\_M}}.
\]
Ошибка высоты:
\[
  e_{h,a}=h_a^\star-h_a.
\]
Положительная ошибка означает, что якорю нужно подняться. В коде она
переводится в throttle-приращение с минимальным ненулевым шагом:
\[
  u_{z,a} =
  \begin{cases}
    0, & |e_{h,a}| \approx 0,\\
    \sgn(e_{h,a})
    \max\left(Z_{\min},
      Z_{\max}\min\left(1,\frac{|e_{h,a}|}{2}\right)\right),
    & \text{иначе},
  \end{cases}
\]
где $Z_{\max}=\text{\code{Z\_PWM\_MAX}}$,
$Z_{\min}=\text{\code{Z\_PWM\_MIN\_STEP}}$.

%====================================================================
\section{Начальное разнесение по высоте}

Если все агенты стартуют на одинаковой высоте, локальный закон может оказаться
симметричным: у всех близкие сигналы, и строй не начинает расходиться. Для
разрушения этой симметрии перед основным контуром выполняется короткий
bootstrap:
\[
  PWM_{\mathrm{thr},i}
  =
  PWM_0+
  \left(B_0+B_1\max(1,i-a)\right),
  \qquad i\ne a,
\]
в течение времени $T_b=\text{\code{BOOTSTRAP\_ALT\_SEPARATION\_SEC}}$.

В текущих константах:
\[
  T_b=2.0\ \text{с},\qquad B_0=90,\qquad B_1=35.
\]

%====================================================================
\section{Итоговый алгоритм одного шага управления}

Для каждого агента $i$ на каждом шаге с периодом
\[
  T_c=\frac{1}{\text{\code{CONTROL\_HZ}}}=0.05\ \text{с}
\]
выполняются операции:
\begin{enumerate}
  \item получить $p_i$ и $v_i$;
  \item получить позиции остальных дронов в общем NED-кадре;
  \item построить множество видимых векторов $\mathcal{N}_i$;
  \item найти ближайшие расстояния по полуосям и заполнить отсутствующие
  горизонтальные значения нулём, а вертикальные значения через
  $R_{\mathrm{vis}}$ и $w$;
  \item вычислить $\sigma_{x,i},\sigma_{y,i},\sigma_{z,i}$;
  \item получить $u_{x,i},u_{y,i}$ из непрерывного соседского sigma-закона
  по $\sigma_x,\sigma_y$ и $u_{z,i}$ из PI по $\sigma_z$;
  \item преобразовать $u_x,u_y$ из $[-1,1]$ и $u_z$ как PWM-приращение в RC PWM:
  roll, pitch, throttle, yaw;
  \item отправить RC override в MAVLink worker.
\end{enumerate}

%====================================================================
\section{Комментарии к смыслу параметров}

\begin{center}
\begin{tabular}{|>{\ttfamily}l|p{9cm}|}
\hline
R\_VIS\_M & Радиус видимости соседей $R_{\mathrm{vis}}$, м. \\
\hline
W\_SPACING\_M & Желаемый вертикальный шаг $w$, м. Должен быть меньше $R_{\mathrm{vis}}$. \\
\hline
AXIS\_EPS\_M & Порог, ниже которого проекция не считается положительной или отрицательной. \\
\hline
XY\_SIGMA\_DEADBAND & Мёртвая зона горизонтальных сигналов $\sigma_x,\sigma_y$. \\
\hline
XY\_PWM\_MAX & Устаревший параметр совместимости CLI; в текущем X/Y-преобразовании не используется. \\
\hline
XY\_PWM\_MIN\_STEP & Устаревший параметр совместимости CLI; в текущем X/Y-преобразовании не используется. \\
\hline
Z\_SIGMA\_DEADBAND & Мёртвая зона вертикального сигнала $\sigma_z$. \\
\hline
Z\_SIGMA\_KP, Z\_SIGMA\_KI & Коэффициенты PI-регулятора вертикального распределения. \\
\hline
Z\_PWM\_MAX & Максимальное PWM-приращение throttle для вертикального PI. \\
\hline
Z\_PWM\_MIN\_STEP & Минимальный ненулевой шаг throttle для удержания высоты якоря. \\
\hline
\end{tabular}
\end{center}

%====================================================================
\section{Краткая интерпретация}

Сценарий можно рассматривать как комбинацию двух задач:
\begin{itemize}
  \item \textbf{горизонтальное распределение}: агенты используют локальные
  расстояния до ближайших видимых соседей по $x,y$ и непрерывные команды
  $\clip(-\sigma_x,-1,1),\clip(-\sigma_y,-1,1)$;
  \item \textbf{вертикальное распределение}: агенты используют локальные
  расстояния до ближайших видимых соседей по $z$ и стремятся сделать эти
  расстояния близкими к $w$.
\end{itemize}

Именно поэтому в коде присутствует общий соседский sigma-закон для всех осей,
но разные исполнительные преобразования: $x,y$ подаются как нормированные
непрерывные команды с линейным масштабированием в RC PWM, а $z$ проходит через PI по $\sigma_z$, адаптированный к
инерционному поведению мультикоптера в режиме \code{POSHOLD}.

\end{document}
